

For the original marked cusp, let
\[
 Z(q)=\frac{\log q}{2\pi i}B_0+C(q),\qquad
 B_0=\begin{pmatrix}0&1\\-1&0\end{pmatrix},
\]
where $C(q)$ is holomorphic and bounded on a disc about zero.
For a sufficiently small closed disc the actual six-dimensional
cusp tube is orientation-preservingly diffeomorphic, relative to
its full marked torus-bundle boundary, to the unitless toric
quotient defined by $C=0$.  The latter has the following finite
normal-crossings plumbing description.

Start with the canonical complex line bundle of
$S=\mathbb {CP}^2\#3\overline{\mathbb {CP}}^{\,2}$ and its six
toric curves in cyclic order
\[
 \begin{aligned}
 C_0&=E_1,& C_1&=H-E_1-E_2,& C_2&=E_2,\\
 C_3&=H-E_2-E_3,& C_4&=E_3,& C_5&=H-E_1-E_3.
 \end{aligned}
 \tag{E.T1}\label{maps:eq:E.T1}
\]
Pair $C_i$ with $C_{i+3}$ for $i=0,1,2$.  Both the surface-normal
line and the canonical line restricted to every $C_i$ have degree
$-1$.  In the resulting $\mathcal O(-1)\oplus\mathcal O(-1)$
coordinates $(z,a,b)$, with $z=0$ at the endpoint meeting
$C_{i+1}$ and $b$ the surface-normal coordinate, the pairing is
\[
             (z,a,b)\longmapsto(z^{-1},zb,za).
 \tag{E.T2}\label{maps:eq:E.T2}
\]
The function giving the cusp parameter is $q=zab$ in this
coordinate chart.  Formula \eqref{maps:eq:E.T2}, not merely the induced map
of zero sections, specifies the normal framing of each plumbing.
The two triple points are the endpoint orbits
$\{0,2,4\}$ and $\{1,3,5\}$, where endpoint $i$ means
$C_i\cap C_{i+1}$.  Ordered local complex triples undergo a
cyclic permutation around each orbit, with threefold composition
the identity.  The plumbing is formed with collars and rounded
corners; its distinguished outer boundary is $|q|=\eta$.

This statement supplies the relative smooth cusp input for the
ordinary framed handle construction of \ref{maps:src:E.2.21}.  It does not
promote the homological cancellations of \ref{maps:src:E.2.16} to smooth
cancellations.

\begin{proof}
We first construct the comparison by a smooth family, rather
than use the logarithmic homeomorphism in \eqref{maps:eq:B.5.7}.  Set
$C_s(q)=(1-s)C(q)$ for $0\leq s\leq1$.  The affine toric cover
is fixed.  On its dense torus the deck translation for
$\lambda\in\mathbb Z^2$ becomes
\[
 (x,q,s)\longmapsto
 \bigl(q^{B_0\lambda}
       \exp(2\pi i C_s(q)\lambda)x,\ q,\ s\bigr).
 \tag{E.T3}\label{maps:eq:E.T3}
\]
The exponential is componentwise.  The powers of $q$ permute
the toric charts by translation of the $A_2$ triangulation,
while the exponential factors are nowhere-zero analytic units
on every chart.  Thus \eqref{maps:eq:E.T3} defines a smooth family of
holomorphic deck actions on a fixed smooth toric covering
space.  On a central stratum the lattice translation permutes
its bounded fan cells without fixed points.  Off the central
fibre, freeness and properness are exactly the nondegeneracy
of the real period lattice.

These assertions hold uniformly after shrinking $\eta$.
Indeed put $T=-\log|q|/(2\pi)$, and let $K$ bound
$\|\operatorname{Im}C_s(q)\|$ on the compact parameter interval
and closed disc.  For $T>K+1$ the matrix
$TB_0+\operatorname{Im}C_s(q)$ is invertible, since $B_0$ is
orthogonal.  A fundamental parallelogram has logarithmic
position
$B_0a+T^{-1}\operatorname{Im}C_s(q)a$ with
$a\in[0,1]^2$; it meets a uniformly bounded number of toric
charts.  More explicitly choose a triangle containing the
unperturbed point $B_0a$ and let $P$ be its unimodular vertex
matrix.  Its barycentric coordinates $t_j$ are nonnegative,
and the corresponding chart coordinates satisfy
$\log|z_j|=t_j\log|q|+
O(2\pi\|P^{-1}\|K\sqrt2)$.
There are only finitely many such triangles, so with
$c=2\pi\sqrt2\max\|P^{-1}\|$ every chosen chart coordinate
has modulus at most $\exp(cK)$, uniformly in $q,s$ and
all compact phases.  Their closed polydiscs have compact
closure in the affine charts.  These finitely many compact
chart images prove
compactness of the quotient over $[0,1]$ and $|q|\leq\eta$.
The central charts and their incidence relations are fixed.
Projection of this quotient to $s$ is consequently a proper
smooth submersion, including on every normal-crossings stratum.
Its only boundary is $|q|=\eta$ together with the parameter
end faces, and its boundary restriction is also a submersion.

On $|q|=\eta$ prescribe the marked real coordinates
$x=\exp(2\pi i(Z_s(q)a+b))$, $(a,b)\in T^4$.
They give a smooth product trivialization in $s$ and precisely
the prescribed mapping-torus marking.  Extend its horizontal
vector field over a collar, and then over the compact family.
Local toric product charts provide extensions tangent to
every central stratum; a partition of unity preserves that
property and the lift of $\partial_s$.  Integrating this
field proves the asserted relative diffeomorphism.  Its
orientation is the complex orientation on each fibre.
This is relative Ehresmann transport in an actual smooth
family; no smoothness is inferred from a logarithmic formula
at $q=0$.

It remains to identify the endpoint $C=0$ with the framed
plumbing data.  In the lattice $\mathbb Z^2$ use cyclic rays
\[
 (1,0),\ (0,1),\ (-1,1),\ (-1,0),\ (0,-1),\ (1,-1).
 \tag{E.T4}\label{maps:eq:E.T4}
\]
The star of the vertex zero in the $A_2$ triangulation consists
of the six triangles spanned by zero and consecutive rays.
The corresponding three-dimensional cones have primitive rays
$(v,1)$ and determinant one.  Their common divisor at the
central ray is the toric surface $S$.  A principal central
fibre has normal line
\[
             \mathcal O_S(-C_0-\cdots-C_5)=K_S.
 \tag{E.T5}\label{maps:eq:E.T5}
\]
The cyclic surface fan in \eqref{maps:eq:E.T4} also proves this identification
directly: it is the fan of the three torus-fixed blow-ups of
$\mathbb {CP}^2$, and each neighboring-ray relation gives
$C_i^2=-1$.  In the Picard basis \eqref{maps:eq:E.T1},
$K_S=-3H+E_1+E_2+E_3$, so $K_S\cdot C_i=-1$.

Near a double curve the two adjacent unimodular charts give
the total space of $\mathcal O(-1)\oplus\mathcal O(-1)$.
Its two base charts have coordinates
\[
 (z,a,b),\qquad
 (w,A,B)=(z^{-1},za,zb),\qquad q=zab=wAB.
 \tag{E.T6}\label{maps:eq:E.T6}
\]
The period translation taking a neighboring component back
to the chosen component exchanges the two normal directions.
It reverses the chosen endpoint coordinate on the opposite
toric curve.  In the first chart this is exactly \eqref{maps:eq:E.T2};
in the target infinity chart it is simply $(z,b,a)$, so it
extends smoothly over the endpoint.  The map is involutive,
preserves $q$, and its complex Jacobian determinant is $1$.
These equalities specify the normal-disk attachment, including
its sign and angular twisting.

The endpoint $C_i\cap C_{i+1}$ is sent to
$C_{i+2}\cap C_{i+3}$.  At the first endpoint write the
ordered ambient coordinates as the surface-normal coordinate,
the $C_i$-normal coordinate and the $C_{i+1}$-normal coordinate.
The next endpoint has their order cyclically permuted.
After three crossings the same coordinates return, not just
the same unframed planes.  Thus the two endpoint orbits give
ordinary triple charts $q=z_0z_1z_2$, with no additional
clutching loop at either triple point.

Equivalently, for a triangle with vertex matrix $P$ whose
columns are $(v_i,1)$, translation by $n\in\mathbb Z^2$
replaces it by $L_nP$, where
\[
 L_n=\begin{pmatrix}1&0&n_1\\0&1&n_2\\0&0&1\end{pmatrix}.
 \tag{E.T7}\label{maps:eq:E.T7}
\]
It acts as the identity in the corresponding translated
unimodular coordinates, since
$(L_nP)^{-1}L_nP=I$.  This verifies the preceding plumbing
and triple identifications in the actual toric charts.
More explicitly, order the star matrix as
$P_i=((0,1),(r_i,1),(r_{i+1},1))$ and reverse its
columns when using the coordinate order $(z,a,b)$.
For every $i$ the matrix of translation by $-r_i$ from
this chart to the opposite following-endpoint chart is
\[
 \begin{pmatrix}-1&0&0\\1&0&1\\1&1&0\end{pmatrix}.
 \tag{E.T9}\label{maps:eq:E.T9}
\]
Its monomials are exactly \eqref{maps:eq:E.T2}.  In normal-first order,
$P_{i+2}^{-1}L_{-r_i}P_i$ is the cyclic permutation
matrix.  Thus the framing and triple-cycle assertions
are checked for all six actual fan identifications,
not imposed on an abstract plumbing diagram.
Taking collars of the component core, three double-curve
annuli and two triple charts, and rounding their outer
corners, gives the asserted finite plumbing.  These are
subcharts of the smooth toric quotient, so the construction
is Hausdorff and the collar identifications are the indicated
smooth ones; it is not a quotient formed by identifying
arbitrary closed disk-bundle neighborhoods.

Finally the period translations for $g,u$ are $-e_2,e_1$.
On $q=\eta e^{-2\pi i\theta}$ the marked expression is
$x=\exp(2\pi i((-\theta+iT)B_0a+b))$.
At $\theta=1$ it agrees with $\theta=0$ and
$b$ replaced by $b-B_0a$.  This is exactly the inverse
endpoint convention for
\[
 M_0=\begin{pmatrix}
 1&0&0&0\\0&1&0&0\\0&1&1&0\\-1&0&0&1
 \end{pmatrix}.
 \tag{E.T8}\label{maps:eq:E.T8}
\]
The comparison therefore preserves the full outer marking,
not merely the isomorphism class of the boundary manifold.
The six unimodular charts, all normal degrees, the plumbing
Jacobian and involution, both triple cycles and \eqref{maps:eq:E.T8} are
checked exactly in \path{code/maps_cusp_smooth_plumbing.py}.
\end{proof}



